\ifdefined\gkver\else\def\gkver{student}\fi
\documentclass[\gkver]{gaokaozhenti}
\begin{document}

\chapter{2012年上海}

\begin{enumerate}
\item
%% number: 1
%% paperName: 2012年上海高考第 1 题 2 分
%% typeId: 1
%% type: 单选题
%% chapter: 光学
%% point: 光电效应
%% method: 无
%% score: 2
%% degree: 850
%% duplicateId: 0
%% body:
在光电效应实验中，用单色光照射某种金属表面，有光电子逸出，则光电子的最大初动能取决于入射光的 \xzanswer{A}
\fourchoices[answer=A]
{频率}
{强度}
{照射时间}
{光子数目}
%% answer: A
%% memo:
\memoanswer{
光电效应的实验现象是：当用频率较高的单色光照射某种金属表面时，发现有光电子逸出金属表面。爱因斯坦用“光子说”成功地解释了光电效应现象：光是由粒子（光子）组成的，每一个光子所具有的能量是h$\nu$，当用某种单色光照射某种金属表面时，光子能量的一部分用来克服金属对电子的束缚即逸出功（或脱出功），其余的能量则全部转化为光电子的动能，所以光电子的最大初动能由光的频率决定。光电效应实验表明，只要光的频率超过金属的红限就有光电子逸出，频率越高则光电子的最大初动能越大，否则光电子的最大初动能就小。光的强度只表明具有该种能量的光子数量，光的强度大则逸出的光电子数多，否则则少。光电效应实验又表明，光电子的逸出是瞬时发生的，与照射时间无关。本题的正确选项是（A）。
}

\item
%% number: 2
%% paperName: 2012年上海高考第 2 题 2 分
%% typeId: 1
%% type: 单选题
%% chapter: 光学
%% point: 光的干涉
%% method: 无
%% score: 2
%% degree: 850
%% duplicateId: 0
%% body:
下图为红光或紫光通过双缝或单缝所呈现的图样，则 \xzanswer{B}
\fourpicture[numstyle=chinese, labelprefix={}, h=1.1cm]
{figs/02a.png}
{figs/02b.png}
{figs/02c.png}
{figs/02d.png}
\fourchoices[answer=B]
{甲为紫光的干涉图样}
{乙为紫光的干涉图样}
{丙为红光的干涉图样}
{丁为红光的干涉图样}
%% answer: B
%% memo:
\memoanswer{
光的干涉现象和衍射现象体现了光具有波动性。光的干涉现象和衍射现象产生的机制不同，它们在观察屏上得到的图样也不同。光的干涉图样通常是由光经过双缝后获得的，最典型的就是杨氏双缝干涉图样；光的干涉图样是明暗相间的条纹，且明条纹或暗条纹乏间的距离总是相等的，而条纹之间距离的大小取决于干涉实验中所使用的单色光的波长，单色光的波长越长，则干涉图样中条纹之间的距离越大，反之则越小。光的衍射图样是光经过单缝后得到的图样；光的衍射图样的特点是，中央明条纹的宽度最大，随着明暗条纹离中央距离的加大，明暗条纹的宽度逐渐变小。由本题提供的四幅图像可以判断出，甲、乙是光的干涉图样，且甲图中光的波长大于乙图中光的波长。丙、丁为光的衍射图样，且丙图中光的波长大于丁图中光的波长。综合起来，显然本题的正确选项是（B），因为红光的波长大于紫光的波长。
}

\item
%% number: 3
%% paperName: 2012年上海高考第 3 题 2 分
%% typeId: 1
%% type: 单选题
%% chapter: 原子和原子核
%% point: 天然放射性
%% method: 无
%% score: 2
%% degree: 850
%% duplicateId: 0
%% body:
与原子核内部变化有关的现象是 \xzanswer{C}
\fourchoices[answer=C]
{电离现象}
{光电效应现象}
{天然放射现象}
{$\alpha$ 粒子散射现象}
%% answer: C
%% memo:
\memoanswer{
原子核内部是有结构的，与此有关的就是天然放射性现象，这也是 19 世纪末物理学家的三大重要发现之一，因此本题的正确选项是（C）。电离现象是指原子核外的电子受激发后可以跃迁到脱离该原子的现象，也就是说当电子获得了足够多的能量后，原来受束缚的电子成为自由电子，即被电离，如氢原子中的电子获得 13.6 eV 的能量后可以被电离，所以电离现象与原子核内部变化无关。光电效应是指光子入射到某些金属表面后可以使金属中的电子逸出的现象，与原子核内部的变化无关。$\alpha$ 粒子散射现象是指 $\alpha$ 粒子射到金箔后被弹回改变了轨迹方向的实验，卢瑟福的 $\alpha$ 粒子散射实验证实了原子中存在原子核。

本题的正确选项为（C）。
}

\item
%% number: 4
%% paperName: 2012年上海高考第 4 题 2 分
%% typeId: 1
%% type: 单选题
%% chapter: 光学
%% point: 光子说
%% method: 无
%% score: 2
%% degree: 850
%% duplicateId: 0
%% body:
根据爱因斯坦的“光子说”可知 \xzanswer{B}
\fourchoices[answer=B]
{“光子说”本质就是牛顿的“微粒说”}
{光的波长越大，光子的能量越小}
{一束单色光的能量可以连续变化}
{只有光子数很多时，光才具有粒子性}
%% answer: B
%% memo:
\memoanswer{
在物理学的发展过程中，牛顿运用弹性小球在分界面上发生碰撞后反弹的理论，提出了光是由弹性小球组成的一束粒子流的思想，用牛顿力学的理论来解释光的反射现象获得成功，由此牛顿提出了光的“微粒说”。虽然牛顿的“微粒说”能解释光的反射和折射，但有局限性。17$\sim$18世纪，人们对光的波动理论有了新的认识。爱因斯坦在研究光电效应时发现用“微粒说”无法解释光电效应现象，由此提出了“光子说”。后来，科学家经过研究发现，光具有波粒二象性，即在某些情况下光表现为粒子性，在另一些情况下光又表现为波动性。显然，“光子说”本质就是牛顿的“微粒说”的说法不成立，选项（A）错误。选项（B）涉及光的波长和能量，因为光的速度为定值，所以波长越大，频率越小，即光子的能量越小，反之则波长越短能量越大，所以本题的正确选项是（B）。既然是单色光，那么能量是确定值，不可能连续变化，选项（C）错误。此外，光的粒子性与光子的多少无关，如果光的频率大于某种金属的红限，那么即使是单个光子也能产生光电效应。选项（D）错误。
}

\item
%% number: 5
%% paperName: 2012年上海高考第 5 题 2 分
%% typeId: 1
%% type: 单选题
%% chapter: 原子和原子核
%% point: 天然放射性
%% method: 无
%% score: 2
%% degree: 650
%% duplicateId: 0
%% body:
在轧制钢板时需要动态地监测钢板厚度，其检测装置由放射源、探测器等构成，如图所示。该装置中探测器接收到的是\xzanswer{D}
\onepicture{\includesvg{figs/05}}
\fourchoices[answer=D]
{X 射线}
{$\alpha$ 射线}
{$\beta$ 射线}
{$\gamma$ 射线}
%% answer: D
%% memo:
\memoanswer{
本题中要利用放射源发出的射线穿透钢板，因此这种射线的穿透性要强。在 $\alpha$、$\beta$、$\gamma$ 这三种射线和X射线中，$\alpha$ 射线和 $\beta$ 射线的穿透能力都很弱，只有 $\gamma$ 射线具有可以穿透钢板的穿透能力；X 射线虽然也有一定的穿透性，但题目中说明，该射线是由放射源发出的，而 X 射线不是由放射源发出的，所以探测器接收到的不可能是 X 射线。

本题的正确选项是（D）。
}

\item
%% number: 6
%% paperName: 2012年上海高考第 6 题 2 分
%% typeId: 1
%% type: 单选题
%% chapter: 力物体的平衡
%% point: 力的分解
%% method: 无
%% score: 2
%% degree: 850
%% duplicateId: 0
%% body:
已知两个共点力的合力为 $50\UN$，分力 $F_{1}$ 的方向与合力 $F$ 的方向成 $30^{\circ}$ 角，分力 $F_{2}$ 的大小为 $30\UN$。则 \xzanswer{C}
\fourchoices[answer=C]
{$F_{1}$ 的大小是唯一的}
{$F_{2}$ 的方向是唯一的}
{$F_{2}$ 有两个可能的方向}
{$F_{2}$ 可取任意方向}
%% answer: C
%% memo:
\memoanswer{
已知 $F_{1}$ 的方向，$F_{2}$ 的最小值为 $F_{2\min}=F\sin30^{\circ}=25\UN$，又因 $25\UN<F_{2}<F=50\UN$，所以 $F_{2}$ 可能有两个可能的方向，如图所示，故 C 正确。
\onepicture[h=2.4cm, align=right]{figs/06_an.jpg}
}

\item
%% number: 7
%% paperName: 2012年上海高考第 7 题 2 分
%% typeId: 1
%% type: 单选题
%% chapter: 电路
%% point: 逻辑电路
%% method: 无
%% score: 2
%% degree: 850
%% duplicateId: 0
%% body:
如图，低电位报警器由两个基本的门电路与蜂鸣器组成，该报警器只有当输入电压过低时蜂鸣器才会发出警报。其中 \xzanswer{B}
\onepicture{\includesvg{figs/07}}
\fourchoices[answer=B]
{甲是“与”门，乙是“非”门}
{甲是“或”门，乙是“非”门}
{甲是“与”门，乙是“或”门}
{甲是“或”门，乙是“与”门}
%% answer: B
%% memo:
\memoanswer{
我们先来分析蜂鸣器报警时的工作情况。蜂鸣器的一端接地，因此只有在乙的输出端为高电位时蜂鸣器才发出警报，因为乙是“非”门，所以要求乙的输入为低电位时蜂鸣器才会报警。如果甲是“与”门，那么甲的输入端 U 无论是低电位还是高电位，它的输出端始终是低电位，也就是说乙的输入始终是低电位，蜂鸣器始终处于报警状态。根据题意，显然这不符合要求。所以甲应该是“或”门，平时输入端 U 为高电位，使乙的输入也为高电位，则乙的输出为低电位，此时蜂鸣器处于非报警状态。而当 U 为低电位时，乙的输入也为低电位，此时乙的输出就为高电位，蜂鸣器就会发出警报。

本题的正确选项是（B）
}

\item
%% number: 8
%% paperName: 2012年上海高考第 8 题 2 分
%% typeId: 1
%% type: 单选题
%% chapter: 运动和力
%% point: 牛顿第二定律
%% method: 整体与隔离
%% score: 2
%% degree: 650
%% duplicateId: 0
%% body:
如图，光滑斜面固定于水平面，滑块 A、B 叠放后一起冲上斜面，且始终保持相对静止，A 上表面水平。则在斜面上运动时，B 受力的示意图为 \xzanswer{A}
\onepicture[h=2.6cm, align=right]{figs/08.png}
\fourchoices[answer=A, ispicture=true, h=2.2cm]
{figs/08a.png}
{figs/08b.png}
{figs/08c.png}
{figs/08d.png}
%% answer: A
%% memo:
\memoanswer{
滑块 A、B 叠放在一起后冲上光滑斜面，且始终保持相对静止，此时滑块 A、B 一起沿斜面做匀减速直线运动，加速度方向沿斜面向下，且滑块 A、B 之间一定有摩擦力作用，所以滑块B应受三个力作用，即重力 $G$、滑块 A 对它的支持力 $F_{N}$ 和摩擦力 $F_{f}$。现在关键是要确定 $F_{f}$ 的方向。因为滑块A上表面光滑，所以 $F_{f}$ 的方向一定沿水平方向，选项（C）、（D）错误。

又因为滑块 B 的加速度沿斜面向下，将此加速度沿水平方向和竖直方向分解可知，滑块 B 一定有一个沿水平方向向左的加速度，也就是说，滑块 B 受到一个水平向左的力作用，此力就是 $F_{f}$，即选项（A）正确，选项（B）错误。

本题的正确选项是（A）。
}

\item
%% number: 9
%% paperName: 2012年上海高考第 9 题 3 分
%% typeId: 1
%% type: 单选题
%% chapter: 原子和原子核
%% point: 原子核的组成
%% method: 无
%% score: 3
%% degree: 850
%% duplicateId: 0
%% body:
某种元素具有多种同位素，反映这些同位素的质量数 $A$ 与中子数 $N$ 关系的是图 \xzanswer{B}
\fourchoices[answer=B, ispicture=true, h=2.4cm]
{figs/09a.png}
{figs/09b.png}
{figs/09c.png}
{figs/09d.png}
%% answer: B
%% memo:
\memoanswer{
图中的横坐标代表中子数，纵坐标表示质量数，质量数是中子数与质子数之和。选项（A）表示中子数在增加，而质量数在减小，这显然是不可能的。选项（B）表示质量数与中子数成正比关系，正确。选项（C）中的 $A$ 与 $N$ 虽然也是正比关系，但是对某种元素的多种同位素来说，它的质子数恒定，而中子数可以有多种不同的数值，所以选项（C）错误。选项（D）表示该元素的中子数在增加，而质量数恒定这也是不可能的。

本题的正确选项是（B）。
}

\item
%% number: 10
%% paperName: 2012年上海高考第 10 题 3 分
%% typeId: 1
%% type: 单选题
%% chapter: 直线运动
%% point: 竖直上下抛运动
%% method: 图象法
%% score: 3
%% degree: 650
%% duplicateId: 0
%% body:
小球每隔 $0.2\Us$ 从同一高度抛出，做初速为 $6\Ums$ 的竖直上抛运动，设它们在空中不相碰。第一个小球在抛出点以上能遇到的小球数为（取 $g=10\Umsq$）\xzanswer{C}
\fourchoices[answer=C]
{三个}
{四个}
{五个}
{六个}
%% answer: C
%% memo:
\memoanswer{
方法一：由 $v_{t}=v_{0}-gt=0$，可得小球上升到最高点需时 $0.6\Us$，即当第三个小球刚抛出时，第一个小球已经到达最高点并开始落下；第一个小球在落下的过程中可以碰到上抛的其他小球，而且当它下落 $0.1\Us$ 时正好与第二个上抛的小球碰上，依次类推，以后每隔 $0.1\Us$ 碰上一个，而第一个小球从最高点下落到抛出点的时间也为 $0.6\Us$，所以总共可能遇到 5 个小球。

方法二：由 $v_{t}=v_{0}-gt=0$，可得每个小球上升到最高点需时 $0.6\Us$，那么从抛出点到回到抛出点需时 $1.2\Us$。由 $h=v_{0}t-\frac{1}{2}gt^{2}$，可得每个小球上升的最大高度 $h_{max}=6\times0.6\Um-\frac{1}{2}\times10\times(0.6)^{2}\Um=1.8\Um$。由此作出上抛的各小球的 $h-t$ 图像如图所示。从图中可以看出，第一个小球在抛出点以上能遇到的小球数为 5 个。

\onepicture[h=3.4cm, align=right]{figs/10_an.png}

本题的正确选项是（C）。
}

\item
%% number: 11
%% paperName: 2012年上海高考第 11 题 3 分
%% typeId: 1
%% type: 单选题
%% chapter: 电场
%% point: 库仑定律
%% method: 无
%% score: 3
%% degree: 650
%% duplicateId: 0
%% body:
A、B、C 三点在同一直线上，$AB:BC=1:2$，B 点位于 A、C 之间，在 B 处固定一电荷量为 $Q$ 的点电荷。当在 A 处放一电荷量为 $+q$ 的点电荷时，它所受到的电场力为 $F$；移去 A 处电荷，在 C 处放一电荷量为 $-2q$ 的点电荷，其所受电场力为 \xzanswer{B}
\fourchoices[answer=B]
{$-\frac{F}{2}$}
{$\frac{F}{2}$}
{$-F$}
{$F$}
%% answer: B
%% memo:
\memoanswer{
根据题意，由库仑定律可得 $F = F_{BA} = k\frac{{Qq}}{{{{(AB)}^2}}}$，$F_{BC} = k\frac{{2Qq}}{{{{(BC)}^2}}} = \frac{F}{2}$，考虑到 $F$ 的方向向左，而 $F_{BC}$ 的方向也向左，所以本题的正确选项为（B）。在本题的解答中，有些考生可能只注意了两个力的比值大小而未注意它们的实际方向，可能会错选（A）。
}

\item
%% number: 12
%% paperName: 2012年上海高考第 12 题 3 分
%% typeId: 1
%% type: 单选题
%% chapter: 曲线运动
%% point: 平抛运动
%% method: 无
%% score: 3
%% degree: 650
%% duplicateId: 0
%% body:
如图，斜面上 a、b、c 三点等距，小球从 a 点正上方 O 点抛出，做初速为 $v_{0}$ 的平抛运动，恰落在 b 点。若小球初速变为 $v$，其落点位于 c，则 \xzanswer{A}
\onepicture{\includesvg{figs/12}}
\fourchoices[answer=A]
{$v_{0} < v$ < 2$v_{0}$}
{$v$ = 2$v_{0}$}
{2$v_{0} < v$ < 3$v_{0}$}
{$v$ > 3$v_{0}$}
%% answer: A
%% memo:
\memoanswer{
\onepicture[h=3.6cm, align=right]{figs/12_an.png}

由题意可知，小球初速变为 $v$ 时，其落点位于 c 点，根据平抛运动的特点，初速度越大，则落点越远，显然 $v_{0} < v$。由于斜面上 a、b、c 三点等距，如图所示，设想作一条过 b 点的水平线，当小球从 a 点抛出的初速度变为 2$v_{0}$ 时，小球恰好在 c 点正下方通过这条水平线上的点 c$'$，然后落到斜面上 c 点下面的点 d，因此可以判定 $v <$ 2$v_{0}$。

本题的正确选项是（A）。
}

\item
%% number: 13
%% paperName: 2012年上海高考第 13 题 3 分
%% typeId: 1
%% type: 单选题
%% chapter: 电路
%% point: 电功电功率
%% method: 无
%% score: 3
%% degree: 850
%% duplicateId: 0
%% body:
当电阻两端加上某一稳定电压时，通过该电阻的电荷量为 $0.3\UC$，消耗的电能为 $0.9\UJ$。为在相同时间内使 $0.6\UC$ 的电荷量通过该电阻，在其两端需加的电压和消耗的电能分别是 \xzanswer{D}
\fourchoices[answer=D]
{$3\UV$，$1.8\UJ$}
{$3\UV$，$3.6\UJ$}
{$6\UV$，$1.8\UJ$}
{$6\UV$，$3.6\UJ$}
%% answer: D
%% memo:
\memoanswer{
本题可用尝试法求解。由 $q=It$，$W=IUt$ 关系，设 $t=1\Us$，则 $I_{1}=0.3\UA$，$U_{1}=W_{1}/I_{1}=3\UV$。在相同时间内（$1\Us$）使 $0.6\UC$ 的电荷量通过该电阻时，$I_{2}=0.6\UA$，即 $W_{2}/U_{2}=0.6\UA$，对照题目所给的选项可知，选项（A）和（D）符合要求，但是选项（A）中的 $3\UV$ 是 $U_{1}$，即加上此电压时，通过该电阻的电荷量为 $0.3\UC$，显然不符合要求。所以本题的正确选项是（D）。

在解这一类试题时，要考虑数值的合理性，如果本解中设 $t=0.1\Us$，则电流为 $I_{1}=3\UA$，这也符合实际，结果仍然是选项（D）正确。所以这种假设是可行的。当然题目的设计也要注意所给出数据的合理性。
}

\item
%% number: 14
%% paperName: 2012年上海高考第 14 题 3 分
%% typeId: 1
%% type: 单选题
%% chapter: 力物体的平衡
%% point: 力矩平衡
%% method: 极值
%% score: 3
%% degree: 650
%% duplicateId: 0
%% body:
如图，竖直轻质悬线上端固定，下端与均质硬棒 AB 中点连接，棒长为线长的二倍。棒的 A 端用铰链墙上，棒处于水平状态。改变悬线的长度，使线与棒的连接点逐渐右移，并保持棒仍处于水平状态。则悬线拉力 \xzanswer{A}
\onepicture{\includesvg{figs/14}}
\fourchoices[answer=A]
{逐渐减小}
{逐渐增大}
{先减小后增大}
{先增大后减小}
%% answer: A
%% memo:
\memoanswer{
\onepicture[h=3.2cm, align=right]{\includesvg{figs/14_an}}

由图可知，线与棒的连接点逐渐右移后，悬线拉力的力臂 $l_{T}$ 逐渐增大，当连接点移至 B 端时力臂最大。由力矩平衡 $M_{G} = T\cdot l_{T}$ 可知，线拉力逐渐减小。
}

\item
%% number: 15
%% paperName: 2012年上海高考第 15 题 3 分
%% typeId: 1
%% type: 单选题
%% chapter: 功和能
%% point: 功
%% method: 无
%% score: 3
%% degree: 650
%% duplicateId: 0
%% body:
质量相等的均质柔软细绳 A、B 平放于水平地面，绳 A 较长。分别捏住两绳中点缓慢提起，直到全部离开地面，两绳中点被提升的高度分别为 $h_{A}$、$h_{B}$，上述过程中克服重力做功分别为 $W_{A}$、$W_{B}$。若 \xzanswer{B}
\fourchoices[answer=B]
{$h_{A} = h_{B}$，则一定有 $W_{A} = W_{B                        }$}
{$h_{A} > h_{B}$，则可能有 $W_{A} < W_{B}$}
{$h_{A} < h_{B}$，则可能有 $W_{A} = W_{B                        }$}
{$h_{A} > h_{B}$，则一定有 $W_{A} > W_{B}$}
%% answer: B
%% memo:
\memoanswer{
\onepicture[h=3.4cm, align=right]{\includesvg{figs/15_an}}

如图所示。当 $h_{A} = h_{B}$ 时，由于绳 A 较长，因此两绳的重心高度 $h_{A}' < h_{B}'$；因为细绳质量相等，由 $W_{G} = mgh$ 可知 $W_{A} < W_{B}$。选项 A 错误。

若 $h_{A} < h_{B}$，必有 $h_{A}' < h_{B}'$，则 $W_{A} < W_{B}$。选项 C 错误。

若 $h_{A} > h_{B}$，则 $h_{A}'$ 和 $h_{B}'$ 的高度的大小关系都有可能。选项 D 错误，B 正确。
}

\item
%% number: 16
%% paperName: 2012年上海高考第 16 题 3 分
%% typeId: 1
%% type: 单选题
%% chapter: 功和能
%% point: 机械能守恒定律
%% method: 整体与隔离
%% score: 3
%% degree: 650
%% duplicateId: 0
%% body:
如图，可视为质点的小球 A、B 用不可伸长的细软轻线连接，跨过固定在地面上半径为 $R$ 的光滑圆柱，A 的质量为 B 的两倍。当 B 位于地面时，A 恰与圆柱轴心等高。将 A 由静止释放，B 上升的最大高度是 \xzanswer{C}
\onepicture{\includesvg{figs/16}}
\fourchoices[answer=C]
{2$R    $}
{$\frac{{5R}}{3}$}
{$\frac{{4R}}{3}$}
{$\frac{{2R}}{3}$}
%% answer: C
%% memo:
\memoanswer{
小球A由静止释放到落到地面时整个系统机械能守恒，则 $m_{A}gR=m_{_{B}}gR+\frac{1}{2}m_{B}v_{B}^{2}+\frac{1}{2}m_{A}v_{A}^{2}$。由于 A、B 小球用不可伸长的细软轻线连接，因此由竖直上抛运动知识可知$\frac{1}{2}m_{B}v_{B}^{2}=m_{B}gh$，$v_{B}=v_{A}$，得：3$v_{B}^{2}=$2$gR$，$v_{B}^{2}=\frac{2}{3}gR$，$v_{B}$ 是 B 球上升到图示圆柱直径处时上抛时的初速度，此后 B 球做竖直上抛运动，B 球的动能全部转化为势能。可求得此后 B 球能上升的高度 $h=\frac{1}{3}R$，加上此前B球已经上升了 $R$，所以 B 球能上升的最大高度为$\frac{4}{3}R$。

本题的正确选项是（C）。
}

\item
%% number: 17
%% paperName: 2012年上海高考第 17 题 4 分
%% typeId: 2
%% type: 多选题
%% chapter: 电路
%% point: 动态分析
%% method: 无
%% score: 4
%% degree: 650
%% duplicateId: 0
%% body:
直流电路如图所示，在滑动变阻器的滑片 P 向右移动时，电源的 \xzanswer{ABC}
\onepicture{\includesvg{figs/17}}
\fourchoices[answer=ABC]
{总功率一定减小}
{效率一定增大}
{内部损耗功率一定减小}
{输出功率一定先增大后减小}
%% answer: ABC
%% memo:
\memoanswer{
当滑片 P 向右滑动时，电阻 $R$ 增大，电流 $I$ 减小，内电压减小，外电压 $U$ 增大，所以电源的总功率 $P = EI$ 减小，内部损耗功率 $I^{2}r$ 减小，电源的效率 $\eta = \frac{{UI}}{{EI}} = \frac{U}{E}$，$U$ 增大，故效率增大。ABC 均正确；

当 $R = r$ 时输出功率最大，由于不知道变阻器原电阻 $R$ 与内电阻 $r$ 的大小关系，所以无法判断电源输出功率的变化情况，D 错误。
}

\item
%% number: 18
%% paperName: 2012年上海高考第 18 题 4 分
%% typeId: 2
%% type: 多选题
%% chapter: 功和能
%% point: 功率
%% method: 无
%% score: 4
%% degree: 450
%% duplicateId: 0
%% body:
位于水平面上的物体在水平恒力 $F_{1}$ 作用下，做速度为 $v_{1}$ 的匀速运动；若作用力变为斜向上的恒力 $F_{2}$，物体做速度为 $v_{2}$ 的匀速运动，且 $F_{1}$ 与 $F_{2}$ 功率相同。则可能有 \xzanswer{BD}
\onepicture{\includesvg{figs/18}}
\fourchoices[answer=BD]
{$F_{2} = F_{1}$，$v_{1} > v_{2                  }$}
{$F_{2} = F_{1}$，$v_{1} < v_{2}$}
{$F_{2} > F_{1}$，$v_{1} > v_{2                   }$}
{$F_{2} < F_{1}$，$v_{1} < v_{2}$}
%% answer: BD
%% memo:
\memoanswer{
由平面上的物体在水平恒力$F_{1}$作用下做速度为$v_{1}$的匀速运动可知，该物体受到一水平向左的摩擦力$f_{1}$作用，摩擦力与$F_{1}$平衡，即$f_{1}=F_{1}$。若作用力变为斜向上的恒力$F_{2}$，物体做速度为$v_{2}$的匀速运动，且$F_{1}$与$F_{2}$功率相同，设此时$F_{2}$与水平面间的夹角为$\theta$，则$F_{1}v_{1}=F_{2}v_{2}$cos$\theta$，$F_{2}$向上的分力作用使物体对水平面的正压力减小，此时物体所受的摩擦力$f_{2}$也比$f_{1}$小，即有$f_{2}<f_{1}$，而$f_{2}=F_{2}$cos$\theta$，即$F_{2}$cos$\theta<F_{1}$。对于上式，有可能$F_{2}>F_{1}$、$F_{2}=F_{1}$和$F_{2}<F_{1}$。

若$F_{2}>F_{1}$，因为此时$F_{2}$cos$\theta<F_{1}$，由$F_{1}v_{1}=F_{2}v_{2}$cos$\theta$得，$v_{1}<v_{2}$，选项（C）错误。

若$F_{2}=F_{1}$，由$F_{1}v_{1}=F_{2}v_{2}$cos$\theta$得，$v_{1}=v_{2}$cos$\theta$，则$v_{1}<v_{2}$，选项（A）错误，选项（B）正确。

若$F_{2}<F_{1}$，由$F_{1}v_{1}=F_{2}v_{2}$cos$\theta$得，$v_{1}<v_{2}$cos$\theta$，则$v_{1}<v_{2}$，选项（D）正确。

即本题的正确选项是（B）、（D）。
}

\item
%% number: 19
%% paperName: 2012年上海高考第 19 题 4 分
%% typeId: 2
%% type: 多选题
%% chapter: 曲线运动
%% point: 圆周运动的应用
%% method: 无
%% score: 4
%% degree: 450
%% duplicateId: 0
%% body:
图 a 为测量分子速率分布的装置示意图。圆筒绕其中心匀速转动，侧面开有狭缝 N，内侧贴有记录薄膜，M 为正对狭缝的位置。从原子炉 R 中射出的银原子蒸汽穿过屏上的 S 缝后进入狭缝 N，在圆筒转动半个周期的时间内相继到达并沉积在薄膜上。展开的薄膜如图 b 所示，NP，PQ 间距相等。则 \xzanswer{AC}
\onepicture{\includesvg{figs/19}}
\fourchoices[answer=AC]
{到达 M 附近的银原子速率较大}
{到达 Q 附近的银原子速率较大}
{位于 PQ 区间的分子百分率大于位于 NP 区间的分子百分率}
{位于 PQ 区间的分子百分率小于位于 NP 区间的分子百分率}
%% answer: AC
%% memo:
\memoanswer{
从原子炉 R 中射出的银原子蒸汽具有不同的速率，它们在穿过屏上的 S 缝后进入狭缝 N，速率大的分子先到达圆筒对侧，此时圆筒只转过很小角度，所以速率大的分子的落点偏向 M，速率略小的分子的落点偏向 N，从展开的薄膜图上可知，PQ 区间的分子较为密集，表明位于 PQ 区间的分子百分率较大，NP 区间上的分子较为稀疏，表明位于 NP 区间的分子百分率较小。因此根据沉积在薄膜上分子的位置和分子的沉积厚度，即可判断出选项（A）、（C）正确。
}

\item
%% number: 20
%% paperName: 2012年上海高考第 20 题 4 分
%% typeId: 2
%% type: 多选题
%% chapter: 功和能
%% point: 动能定理
%% method: 无
%% score: 4
%% degree: 450
%% duplicateId: 0
%% body:
如图，质量分别为 $m_{A}$ 和 $m_{B}$ 的两小球带有同种电荷，电荷量分别为 $q_{A}$ 和 $q_{B}$，用绝缘细线悬挂在天花板上。平衡时，两小球恰处于同一水平位置，细线与竖直方向间夹角分别为 $\theta_{1}$ 与 $\theta_{2}$（$\theta_{1} > \theta_{2}$）。两小球突然失去各自所带电荷后开始摆动，最大速度分别为 $v_{A}$ 和 $v_{B}$，最大动能分别为 $E_{kA}$ 和 $E_{kB}$。则 \xzanswer{ACD}
\onepicture{\includesvg{figs/20}}
\fourchoices[answer=ACD]
{$m_{A}$ 一定小于 $m_{B                   }$}
{$q_{A}$ 一定大于 $q_{B}$}
{$v_{A}$ 一定大于 $v_{B                      }$}
{$E_{kA}$ 一定大于 $E_{kB}$}
%% answer: ACD
%% memo:
\memoanswer{
\onepicture[h=3.6cm, align=right]{figs/20_an.png}

开始时，A、B 两球在库仑力和重力的作用下处于平衡状态，设库仑力为 $F$，则由 $F$/$m_{A}g =$tan$\theta_{1}$、$F$/$m_{B}g =$tan$\theta_{2}$ 及 $\theta_{1}>\theta_{2}$，得到 $m_{B}>m_{A}$，选项（A）正确。

而$\theta_{1}$与$\theta_{2}$的大小和两个电荷所带电荷量的多少无关，所以选项（B）错误。

令图中竖直虚线的高度为$h$，两侧小球由于库仑力的作用而使悬线偏离竖直方向的角度为$\theta$，则悬线的长度$L=h/\cos\theta$ ，则两侧小球在库仑力作用下摆起的高度为$\Delta h=L(1-\cos\theta)=\frac{h}{\cos \theta }-h$。而当两侧小球的电荷量失去后，两小球向下摆动时遵守机械能守恒定律，因此$mg\Delta h=\frac{1}{2}mv_{max}^{2}$，解得$v_{max}=\sqrt {2gh(\frac{1}{{\cos \theta }} - 1)} $，而$\theta$↑$\to$cos$\theta$↓$\to\frac{1}{{\cos \theta }}$↑$\to v_{max}$↑，由于$\theta_{1}>\theta_{2}$，所以选项（C）正确。

又因为$E_{k}=\frac{1}{2}mv_{max}^{2}=mg\Delta h=F$cot$\theta$（$\frac{h}{{\cos \theta }} -h$）$=Fh$tan$\frac{\theta }{2}$，而$\theta$↑$\to$tan$\frac{\theta }{2}$ ↑$\to E_{k}$↑，同样由于$\theta_{1}>\theta_{2}$，所以选项（D）正确。

本题的正确选项是（A）、（C）、（D）。
}

\item
%% number: 21
%% paperName: 2012年上海高考第 21 题 4 分
%% typeId: 3
%% type: 填空题
%% chapter: 原子和原子核
%% point: 天然放射性
%% method: 无
%% score: 4
%% degree: 850
%% duplicateId: 0
%% body:
$\ce{^{60}_{27}Co}$ 发生一次 $\beta$ 衰变后变为 $\ce{Ni}$，其衰变方程为\_\_\_\_\_\_\_\_\_\_\_，在该衰变过程中还发出频率为 $\nu_{1}$、$\nu_{2}$ 的两个光子，其总能量为 \_\_\_\_\_\_\_\_\_\_\_。
%% answer: 
\jdanswer{
$\ce{^{60}_{27}Co -> ^{60}_{28}Ni + ^{0}_{-1}e}$，$h$（$\nu_{1} + \nu_{2}$）
}
%% memo:
\memoanswer{
$\ce{^{60}_{27}Co}$ 发生一次 $\beta$ 衰变不改变质量数，所以衰变方程为 $\ce{^{60}_{27}Co -> ^{60}_{28}Ni + ^{0}_{-1}e}$；在该衰变过程中还发出频率为 $\nu_{1}$、$\nu_{2}$ 的两个光子，光子的能量为 $h\nu$，其中 $\nu$ 为光子的频率，所以其总能量为 $h$（$\nu_{1}+\nu_{2}$）。

本题的正确答案为：$\ce{^{60}_{27}Co -> ^{60}_{28}Ni + ^{0}_{-1}e}$，$h$（$\nu_{1}+\nu_{2}$）。
}

\item
%% number: 22
%% paperName: 2012年上海高考第 22 题 4 分
%% typeId: 3
%% type: 填空题
%% chapter: 曲线运动
%% point: 人造地球卫星
%% method: 无
%% score: 4
%% degree: 650
%% duplicateId: 0
%% body:
本题为分叉题，分 A、B 两类，考生可任选一类答题，若两类试题均做，一律按 A 类题计分。

\textbf{A．}A、B 两物体在光滑水平地面上沿一直线相向而行，A 质量为 $5\Ukg$，速度大小为 $10\Ums$，B 质量为 $2\Ukg$，速度大小为 $5\Ums$，它们的总动量大小为\_\_\_\_\_\_\_\_\_$\Ukgms$；两者相碰后，A 沿原方向运动，速度大小为 $4\Ums$，则 B 的速度大小为\_\_\_\_\_\_\_\_\_$\Ums$。

\textbf{B．}人造地球卫星做半径为 $r$，线速度大小为 $v$ 的匀速圆周运动。当其角速度变为原来的 $\frac{{\sqrt 2 }}{4}$ 倍后，运动半径为\_\_\_\_\_\_\_\_\_，线速度大小为 \_\_\_\_\_\_\_\_\_。
%% answer: A．40；10  B．2r；$\frac{{\sqrt 2 }}{2}$v
\jdanswer{
A．40；10

B．$2r$；$\frac{{\sqrt 2 }}{2}$v
}
%% memo:
\memoanswer{
\textbf{A．}A、B 两物体在光滑水平地面上沿一直线相向而行过程中发生碰撞，根据动量守恒定律，碰撞前两个物体的初动量之和等于碰撞后的总动量，即 $m_{A}v_{A0}-m_{B}v_{B0}=m_{A}v_{A}+m_{B}v_{B}$。因动量是矢量，在写出动量守恒表达式时要考虑动量的方向（在上式中设向右为正方向），因考虑到 A 物体的质量大于 B 物体的质量，所以 A 物体在碰撞后沿原来的方向运动，B 物体的运动方向应该与 A 物体的相同，由此可求得碰撞前它们的总动量大小为 $40\Ukgms$，碰撞后 B 物体的速度大小为 $10\Ums$。

\textbf{B．}人造地球卫星是受万有引力作用而做圆周运动的，有：$G\frac{{Mm}}{{{r^2}}}=m\frac{{{v^2}}}{r}$，其中 M 为地球质量，m 为卫星质量。当其角速度变为原来的 $\frac{{\sqrt 2 }}{4}$ 倍后，圆周运动的轨道半径 r 将发生变化，由前面的万有引力提供卫星做圆周运动的公式可得：$v^{2}=G\frac{{Mm}}{r}$，根据前后两次的轨道变化，得到 $v_{1}^{2}r_{1}=v_{2}^{2}r_{2}$，将 $v=r\omega$ 代入，可求得 $r_{2}=2r_{1}=2r$。再求线速度，$v_{2}=r_{2}\omega_{2}=2r\cdot\frac{{\sqrt 2 }}{4}\omega_{1}=\frac{{\sqrt 2 }}{2}v$。

本题的正确答案为：A．40；10；B．2r；$\frac{{\sqrt 2 }}{2}$v。
}

\item
%% number: 23
%% paperName: 2012年上海高考第 23 题 4 分
%% typeId: 3
%% type: 填空题
%% chapter: 直线运动
%% point: 运动图象
%% method: 无
%% score: 4
%% degree: 650
%% duplicateId: 0
%% body:
质点做直线运动，其 $s-t$ 关系如图所示，质点在 $0\sim20\Us$ 内的平均速度大小为\_\_\_\_\_\_$\Ums$；质点在\_\_\_\_\_\_\_\_\_时的瞬时速度等于它在 $6\sim20\Us$ 内的平均速度。
\onepicture[align=right]{\includesvg{figs/23}}
%% answer: 
\jdanswer{
0.8；$10\Us$ 或 $14\Us$
}
%% memo:
\memoanswer{
由 $s-t$ 关系图可知，该质点做变速直线运动，质点在 $0\sim20\Us$ 时间内的位移是 $16\Um$，则计算得到的平均速度大小为 $\frac{16}{20}\Ums=0.8\Ums$。求质点在 $6\sim20\Us$ 内的平均速度时，由图可知，在 $6\sim20\Us$ 这 $14\Us$ 内质点的位移为 $15\Um$，即平均速度大小为 $\frac{15}{14}\Ums$。而在 $s-t$ 关系图上，曲线上某点处的斜率表示该时刻的瞬时速度，因此，用直尺过曲线上 $6\Us$ 和 $20\Us$ 的点，平行移动直尺，可发现，直尺在 $10\Us$ 和 $14\Us$ 处与曲线相切，所以质点在 $10\Us$ 或 $14\Us$ 时的瞬时速度等于它在 $6\sim20\Us$ 内的平均速度。

本题的正确答案为：0.8；$10\Us$ 或 $14\Us$。
}

\item
%% number: 24
%% paperName: 2012年上海高考第 24 题 4 分
%% typeId: 3
%% type: 填空题
%% chapter: 机械波
%% point: 波的多解性
%% method: 无
%% score: 4
%% degree: 450
%% duplicateId: 0
%% body:
如图，简谐横波在 $t$ 时刻的波形如实线所示，经过 $\Delta t=3\Us$，其波形如虚线所示。已知图中 $x_{1}$ 与 $x_{2}$ 相距 $1\Um$，波的周期为 $T$，且 $2T<\Delta t<4T$。则可能的最小波速为\_\_\_\_\_\_$\Ums$，最小周期为\_\_\_\_\_\_$\Us$。
\onepicture[align=right]{\includesvg{figs/24}}
%% answer: 
\jdanswer{
5；$\frac{7}{9}$
}
%% memo:
\memoanswer{
由图可知，该波的波长为 $\lambda=7\Um$。而 $x_{1}$ 与 $x_{2}$ 之间的距离为 $1\Um=\frac{1}{7}\lambda$。因为波从实线波形传播到虚线波形用时 $\Delta t=3\Us$，而传播方向未知，所以需讨论。若波向右传播，因为 $2T<\Delta t<4T$，所以从实线波形传播到虚线波形所传播的距离可能为 $2\lambda+\frac{1}{7}\lambda=\frac{15}{7}\lambda$ 或 $3\lambda+\frac{1}{7}\lambda=\frac{22}{7}\lambda$；若波向左传播，同样因为 $2T<\Delta t<4T$，所以从实线波形传播到虚线波形所传播的距离可能为 $3\lambda-\frac{1}{7}\lambda=\frac{20}{7}\lambda$ 或 $4\lambda-\frac{1}{7}\lambda=\frac{27}{7}\lambda$。当波速最小时，波在 $\Delta t=3\Us$ 内传播的距离最小，即为 $\frac{15}{7}\lambda$，所以此时波速应为 $v_{min}=\frac{15}{7}\times7\Um\div3\Us=5\Ums$。又因为 $\Delta t=3\Us$ 可能为 $\frac{15}{7}T$、$\frac{22}{7}T$、$\frac{20}{7}T$ 和 $\frac{27}{7}T$，所以最小周期的求法为 $\frac{27}{7}T_{min}=3\Us$，$T_{min}=\frac{7}{9}\Us$。

本题答案为：5；$\frac{7}{9}$。
}

\item
%% number: 25
%% paperName: 2012年上海高考第 25 题 4 分
%% typeId: 3
%% type: 填空题
%% chapter: 电磁感应
%% point: 电磁感应电路分析
%% method: 无
%% score: 4
%% degree: 650
%% duplicateId: 0
%% body:
正方形导线框处于匀强磁场中，磁场方向垂直框平面，磁感应强度随时间均匀增加，变化率为 $k$。导体框质量为 $m$、边长为 $L$，总电阻为 $R$，在恒定外力 $F$ 作用下由静止开始运动。导体框在磁场中的加速度大小为\_\_\_\_\_\_\_，导体框中感应电流做功的功率为 \_\_\_\_\_\_\_\_\_。
\onepicture[align=right]{\includesvg{figs/25}}
%% answer: 
\jdanswer{
$\frac{F}{m}$；$\frac{{{k^2}{L^4}}}{R}$
}
%% memo:
\memoanswer{
由题意可知，线框在磁场中，因此，线框有两条对边都在切割磁感线，无动生电动势，线框中的感应电动势来自磁场的变化。感应电动势 $E = \frac{{\Delta \Phi }}{{\Delta t}} = kL^{2}$，感应电流 $I = \frac{{k{L^2}}}{R}$。又由于左右两条对边中的感应电流方向相反、大小相等，所以线框所受安培力的合力为零。线框所受的合力为 $F$，由 $F = ma$ 得，$a = \frac{F}{m}$；感应电流做功的功率为 $P = \frac{{{E^2}}}{R} = \frac{{{k^2}{L^4}}}{R}$ 。

本题的答案为：$\frac{F}{m}$；$\frac{{{k^2}{L^4}}}{R}$ 。
}

\item
%% number: 26
%% paperName: 2012年上海高考第 26 题 4 分
%% typeId: 4
%% type: 实验题
%% chapter: 电磁感应
%% point: 验证电磁感应实验
%% method: 无
%% score: 4
%% degree: 650
%% duplicateId: 0
%% body:
为判断线圈绕向，可将灵敏电流计 G 与线圈 L 连接，如图所示。已知线圈由 a 端开始绕至 b 端；当电流从电流计 G 左端流入时，指针向左偏转。
\begin{enumerate}
	\item
	将磁铁 N 极向下从线圈上方竖直插入 L 时，发现指针向左偏转。俯视线圈，其绕向为\_\_\_\_\_\_\_\_\_\_（填“顺时针”或“逆时针 ”）。
	\item
	当条形磁铁从图中虚线位置向右远离 L 时，指针向右偏转。俯视线圈，其绕向为\_\_\_\_\_\_\_\_\_\_（填“顺时针”或“逆时针”）。
\end{enumerate}
\onepicture[align=right]{\includesvg{figs/26}}
%% answer: 
\jdanswer{
（1）顺时针；（2）逆时针。
}
%% memo:
\memoanswer{
（1）当磁铁 N 极插入线圈时，在线圈中将产生感应电动势，该电动势的感应电流产生的磁通量阻碍原磁通量的增加，即线圈的上端为 N 极。而当磁铁 N 极向下从线圈上方竖直插入 L 时，发现电流计指针向左偏转，说明感应电流是从 a 端流出并流入灵敏电流计 G 的，又知道线圈由 a 端开始绕至 b 端，所以俯视线圈，线圈的绕向应为顺时针。

（2）当条形磁铁从图中虚线位置向右远离 L 时，从线圈下面穿入线圈的磁通量在减少，因此线圈中产生的感应电流的磁通量将补偿从线圈下面穿入线圈的磁通量的减少，即线圈的上端为 N 极。而从题意可知，电流计的指针向右偏转，说明感应电流是从 b 端流出并流入电流计的，又知道线圈由 a 端开始绕至 b 端，所以俯视线圈，可知其绕向为逆时针方向。

\twopicture[showlabel=false, h=3.2cm, align=right]{figs/26_an1.png}{figs/26_an2.png}

本题的正确答案为：（1）顺时针；（2）逆时针。
}

\item
%% number: 27
%% paperName: 2012年上海高考第 27 题 6 分
%% typeId: 4
%% type: 实验题
%% chapter: 电路
%% point: 多用表的使用
%% method: 无
%% score: 6
%% degree: 650
%% duplicateId: 0
%% body:
在练习使用多用表的实验中：
\begin{enumerate}
	\item
	某同学连接的电路如图所示

	①若旋转选择开关，使尖端对准直流电流挡，此时测得的是通过\_\_\_\_\_\_\_\_的电流；

	②若断开电路中的电键，旋转选择开关使其尖端对准欧姆挡，此时测得的是\_\_\_\_\_\_\_\_的电阻；

	③若旋转选择开关，使尖端对准直流电压挡，闭合电键，并将滑动变阻器的滑片移至最左端，此时测得的是\_\_\_\_\_\_\_\_两端的电压。

	\item
	（单选）在使用多用表的欧姆挡测量电阻时，若（      ）
	\fourchoices[answer=D]
	{双手捏住两表笔金属杆，测量值将偏大}
	{测量时发现指针偏离中央刻度过大，则必需减小倍率，重新调零后再进行测量}
	{选择“$\times$10”倍率测量时发现指针位于 20 与 30 正中间，则测量值小于 25 $\Omega$}
	{欧姆表内的电池使用时间太长，虽然完成调零，但测量值将略偏大}
\end{enumerate}
\onepicture[align=right]{\includesvg{figs/27}}
%% answer: 
\jdanswer{
\begin{enumerate}
	\item
	①$R_{1}$，②$R_{1}$ 和 $R_{2}$ 串联，③$R_{2}$（或：电源）

	\item
	D

\end{enumerate}
}
%% memo:
\memoanswer{
（1）①旋转选择开关，使尖端对准直流电流挡，此时测得的是通过R$_{1}$的电流；②若断开电路中的电键，此时 R$_{1}$ 和 R$_{2}$ 串联，旋转选择开关使其尖端对准欧姆挡，此时测得的是R$_{1}$和R$_{2}$串联的电阻；③若旋转选择开关，使尖端对准直流电压挡，闭合电键，并将滑动变阻器的滑片移至最左端，此时测得的是 R$_{2}$（或电源）两端的电压。

（2）在使用多用表的欧姆挡测量电阻时，若双手捏住两表笔金属杆，相当于给被测电阻并联了一个电阻，则测量值将偏小，且这是不允许的，所以选项（A）错误。测量时发现指针偏离中央刻度过大，说明选择的倍率有可能太小，也有可能太大，所以可能需要减小倍率，也可能需要增大倍率，重新调零后再进行测量，所以选项（B）错误。选择“$\times$10”倍率测量时发现指针位于 20 与 30 正中间，则测量值在 200 $\Omega$ 与 300 $\Omega$ 之间，显然读数错误，选项（C）错误。当欧姆表内的电池使用时间太长后，电池的内阻增加，叠加在外测电阻上，会使测量值偏大，选项（D）正确。

本题对没有使用过多用表的考生有一定难度，因为这些现象都是在实验中经常碰到的。本题对引导考生重视学生实验有积极意义。
}

\item
%% number: 28
%% paperName: 2012年上海高考第 28 题 6 分
%% typeId: 4
%% type: 实验题
%% chapter: 气体性质
%% point: 气体图象
%% method: 无
%% score: 6
%% degree: 650
%% duplicateId: 0
%% body:
右图为“研究一定质量气体在压强不变的条件下，体积变化与温度变化关系”的实验装置示意图。粗细均匀的弯曲玻璃管 A 臂插入烧瓶，B 臂与玻璃管 C 下部用橡胶管连接，C 管开口向上，一定质量的气体被封闭于烧瓶内。开始时，B、C 内的水银面等高。
\onepicture[h=4cm, align=right]{\includesvg{figs/28}}
\begin{enumerate}
	\item
	若气体温度升高，为使瓶内气体的压强不变，应将 C 管\_\_\_\_\_\_\_（填“向上”或“向下”）移动，直至\_\_\_\_\_\_\_\_\_\_\_\_\_。
	\item
	（单选）实验中多次改变气体温度，用 $\Delta t$ 表示气体升高的温度，用 $\Delta h$ 表示 B 管内水银面高度的改变量。根据测量数据作出的图线是
	\fourchoices[answer=A, ispicture=true, h=2.2cm]
	{figs/28a.png}
	{figs/28b.png}
	{figs/28c.png}
	{figs/28d.png}
\end{enumerate}
%% answer: 
\jdanswer{
\begin{enumerate}
	\item
	向下；B、C 两管内水银面等高

	\item
	A

\end{enumerate}
}
%% memo:
\memoanswer{
（1）当气体温度升高时，瓶内气体的压强将增加，为使瓶内气体的压强不变，只有增大瓶内气体的体积，即应将 C 管向下移动，直至 B、C 两管内水银面等高。若 B 管内水银面高于 C 管内水银面，则说明瓶内气体压强增大，反之则减小，即这两种情况都不能保证瓶内气体的压强不变。

（2）气体温度升高时，B 管内水银面高度的改变量 $\Delta h$ 与温度改变量 $\Delta t$ 成正比关系，所以选项（A）正确，选项（B）错误。选项（C）、（D）表示 $\Delta h$ 与 $\Delta t$ 不是同步变化，这不符合实际，所以都错误。

本题的正确答案为：（1）向下；B、C 两管内水银面等高；（2）A。
}

\item
%% number: 29
%% paperName: 2012年上海高考第 29 题 8 分
%% typeId: 4
%% type: 实验题
%% chapter: 机械振动
%% point: 单摆测重力加速度
%% method: 无
%% score: 8
%% degree: 650
%% duplicateId: 0
%% body:
在“利用单摆测重力加速度”的实验中：
\begin{enumerate}
	\item
	某同学尝试用 DIS 测量周期。如图，用一个磁性小球代替原先的摆球，在单摆下方放置一个磁传感器，其轴线恰好位于单摆悬挂点正下方。图中磁传感器的引出端A应接到\_\_\_\_\_\_\_\_\_\_。使单摆做小角度摆动，当磁感应强度测量值最大时，磁性小球位于\_\_\_\_\_\_\_\_\_\_。若测得连续 $N$ 个磁感应强度最大值之间的时间间隔为 $t$，则单摆周期的测量值为\_\_\_\_\_\_\_\_\_\_（地磁场和磁传感器的影响可忽略）。
	\item
	多次改变摆长使单摆做小角度摆动，测量摆长 $L$ 及相应的周期$T$后，分别取 $L$ 和 $T$ 的对数，所得到的 lg$T$-lg$L$ 图线为\_\_\_\_\_\_（填“直线”、“对数曲线”或“指数曲线”）；读得图线与纵轴交点的纵坐标为 $c$，由此得到该地的重力加速度 $g =$ \_\_\_\_\_\_\_\_\_\_。
\end{enumerate}
\onepicture[align=right]{\includesvg{figs/29}}
%% answer: 
\jdanswer{
\begin{enumerate}
	\item
	数据采集器，最低点（或平衡位置），$\frac{{2t}}{{N - 1}}$

	\item
	直线，$\frac{{4{\pi ^2}}}{{{{10}^{2c}}}}$

\end{enumerate}
}
%% memo:
\memoanswer{
（1）DIS实验仪器主要是由传感器、数据采集器和计算机组成的，传感器测得的信号由数据采集器处理后传给计算机作进一步分析，三者缺一不可。因此，磁传感器的引出端A应接到数据采集器，当磁感应强度测量值最大时，磁性小球位于最低点（或平衡位置）。连续测得 $N$ 个磁感应强度最大值之间的时间间隔为 $t$，则单摆周期的测量值为 $\frac{{2t}}{{N - 1}}$ 。

（2）在实验数据处理中，人们习惯于将所得数据之间的关系用线性关系来描述，对于那些不等距的变化量，通常采用取对数或指数的方法，尽量使数据线性化，所以所得到的 lg$T$-lg$L$ 图线为直线。读得图线与纵轴交点的纵坐标为 $c$，即表明 lg$T=c$时lg$L=$0，可解得 $T=$10$^{c}$时$L=$10$^{0}=$1，代入 $T=$2$\pi\sqrt {\frac{l}{g}} $，得 $g=\frac{{4{\pi ^2}}}{{{{10}^{2c}}}}$。

本题的正确答案为：（1）数据采集器，最低点（或平衡位置），$\frac{{2t}}{{N - 1}}$，（2）直线，$\frac{{4{\pi ^2}}}{{{{10}^{2c}}}}$。
}

\item
%% number: 30
%% paperName: 2012年上海高考第 30 题 10 分
%% typeId: 6
%% type: 计算题
%% chapter: 运动和力
%% point: 牛顿第二定律
%% method: 分情况讨论
%% score: 10
%% degree: 650
%% duplicateId: 0
%% body:
如图，将质量 $m=0.1\Ukg$ 的圆环套在固定的水平直杆上。环的直径略大于杆的截面直径。环与杆间动摩擦因数 $\mu=0.8$。对环施加一位于竖直平面内斜向上，与杆夹角 $\theta=53^{\circ}$ 的拉力 $F$，使圆环以 $a=4.4\Umsq$ 的加速度沿杆运动，求 $F$ 的大小。（取 $\sin53^{\circ}=0.8$，$\cos53^{\circ}=0.6$，$g=10\Umsq$）。
\onepicture[align=right]{\includesvg{figs/30}}
%% answer: 
\jdanswer{
$F=1\UN$ 或 $9\UN$
}
%% memo:
\memoanswer{
\onepicture[h=3.2cm, align=right]{\includesvg{figs/30_an1}}

令 $F\sin53^{\circ} = mg$，$F = 1.25\UN$

当 $F < 1.25\UN$ 时，环与杆的上部接触，受力如图。由牛顿定律

$F\cos\theta - \mu F_{N} = ma$

$F_{N} + F\sin\theta = mg$

由此得

$F = \frac{m(a + \mu g)}{\cos \theta  + \mu \sin \theta } = \frac{0.1 \times (4.4 + 0.8 \times 10)}{0.6 + 0.8 \times 0.8}\UN = 1\UN$

\onepicture[h=3.2cm, align=right]{\includesvg{figs/30_an2}}

当 $F > 1.25\UN$ 时，环与杆的下部接触，受力如图。由牛顿定律

$F\cos\theta - \mu F_{N} = ma$

$F\sin\theta = mg + F_{N}$

由此得

$F = \frac{m(a - \mu g)}{\cos \theta  - \mu \sin \theta } = \frac{0.1 \times (4.4 - 0.8 \times 10)}{0.6 - 0.8 \times 0.8}\UN = 9\UN$
}

\item
%% number: 31
%% paperName: 2012年上海高考第 31 题 12 分
%% typeId: 6
%% type: 计算题
%% chapter: 气体性质
%% point: 玻意耳定律
%% method: 无
%% score: 12
%% degree: 650
%% duplicateId: 0
%% body:
如图，长 $L = 100\Ucm$，粗细均匀的玻璃管一端封闭。水平放置时，长 $L_{0} = 50\Ucm$ 的空气柱被水银柱封住，水银柱长 $h = 30\Ucm$。将玻璃管缓慢地转到开口向下和竖直位置，然后竖直插入水银槽，插入后有 $\Delta h = 15\Ucm$ 的水银柱进入玻璃管。设整个过程中温度始终保持不变，大气压强 $p_{0} = 75\UcmHg$。求：
\begin{enumerate}
	\item
	插入水银槽后管内气体的压强 $p$；
	\item
	管口距水银槽液面的距离 $H$。
\end{enumerate}
\onepicture[align=right]{\includesvg{figs/31}}
%% answer: 
\jdanswer{
\begin{enumerate}
	\item
	$62.5\UcmHg$

	\item
	$27.5\Ucm$

\end{enumerate}
}
%% memo:
\memoanswer{
（1）设当转到竖直位置时，水银恰好位于管口位置而未从管中流出，管截面积为 $S$。此时气柱长度 $l=70\Ucm$。由玻意耳定律

$p=\frac{{{p_0}{L_0}}}{l}=\frac{{75\times50}}{{70}}\UcmHg=53.6\UcmHg$

由于 $p+\rho gh=83.6\UcmHg$ 大于 $p_{0}$，因此必有水银从管中漏出。

设当管转至竖直位置时，管内水银柱长度为 $x$，由玻意耳定律得

$p_{0}SL_{0}=(p_{0}-\rho gx)S(L-x)$

整理并代入数据后得 $75\times50=(75-x)(100-x)$

解得：$x=25\Ucm$

设插入槽内后管内气柱长为 $L'$，由题设条件得 $L'=L-(x+\Delta h)=60\Ucm$

由玻意耳定律，插入后管内压强 $p=\frac{{{p_0}{L_0}}}{{L'}}=\frac{{75\times50}}{{60}}\UcmHg=62.5\UcmHg$

（2）设管内外水银面高度差为 $h'$

$h'=75-62.5=12.5\Ucm$

管口距槽内水银面距离

$H=L-L'-h'=(40-12.5)\Ucm=27.5\Ucm$
}

\item
%% number: 32
%% paperName: 2012年上海高考第 32 题 13 分
%% typeId: 1
%% type: 单选题
%% chapter: 磁场
%% point: 安培力
%% method: 无
%% score: 13
%% degree: 450
%% duplicateId: 0
%% body:
载流长直导线周围磁场的磁感应强度大小为 $B = k\frac{I}{r}$，式中常量 $k >$ 0，$I$ 为电流强度，$r$ 为距导线的距离。在水平长直导线 MN 正下方，矩形线圈 abcd 通以逆时针方向的恒定电流，被两根轻质绝缘细线静止地悬挂，如图所示。开始时 MN 内不通电流，此时两细线内的张力均为 $T_{0}$。当MN通以强度为 $I_{1}$ 的电流时，两细线内的张力均减小为 $T_{1}$，当 MN 内电流强度变为 $I_{2}$ 时，两细线内的张力均大于 $T_{0}$。

（1）分别指出强度为 $I_{1}$、$I_{2}$ 的电流的方向；

（2）求 MN 分别通以强度为 $I_{1}$、$I_{2}$ 的电流时，线框受到的安培力 $F_{1}$ 与 $F_{2}$ 大小之比；

（3）当 MN 内的电流强度为 $I_{3}$ 时两细线恰好断裂，在此瞬间线圈的加速度大小为 $a$，求 $I_{3}$ 。
\onepicture[h=4.5cm, align=right]{\includesvg{figs/32}}
%% answer: IIFFIIII
\jdanswer{
\begin{enumerate}
	\item
	$I_{1}$ 方向向左，$I_{2}$ 方向向右

	\item
	$\frac{F_{1}}{F_{2}} = \frac{I_{1}}{I_{2}}$

	\item
	$I_{3} = \frac{(a-g)T_{0}}{(T_{0}-T_{1})g}I_{1}$

\end{enumerate}
}
%% memo:
\memoanswer{
（1）$I_{1}$ 方向向左，$I_{2}$ 方向向右

（2）当 MN 中通以电流 $I$ 时，线圈所受安培力大小为

$F=kIiL$（$\frac{1}{{{r_1}}} - \frac{1}{{{r_2}}}$）

式中 $r_{1}$、$r_{2}$ 分别为 ab、cd 与 MN 的间距，$i$ 为线圈中的电流，$L$ 为 ab、cd 的长度。

$\frac{{{F_1}}}{{{F_2}}} = \frac{{{I_1}}}{{{I_2}}}$

（3）设 MN 中电流强度为 $I_{3}$ 时，线框所受安培力为 $F_{3}$。

由题设条件有

2$T_{0} = G\text{，}$2$T_{1} + F_{1} = G$，$F_{3} + G = \frac{{2{T_0}}}{g}$a

$\frac{{{I_1}}}{{{I_3}}} = \frac{{{F_1}}}{{{F_3}}} = \frac{{({T_0} - {T_1})g}}{{(a - g){T_0}}}$

${I_3} = \frac{{(a - g){T_0}}}{{({T_0} - {T_1})g}}{I_1}$
}

\item
%% number: 33
%% paperName: 2012年上海高考第 33 题 14 分
%% typeId: 6
%% type: 计算题
%% chapter: 电磁感应
%% point: 电磁感应中的匀变速
%% method: 无
%% score: 14
%% degree: 450
%% duplicateId: 0
%% body:
如图，质量为 $M$ 的足够长金属导轨 abcd 放在光滑的绝缘水平面上。一电阻不计，质量为 $m$ 的导体棒 PQ 放置在导轨上，始终与导轨接触良好，PQbc 构成矩形。棒与导轨间动摩擦因数为 $\mu$，棒左侧有两个固定于水平面的立柱。导轨 bc 段长为 $L$，开始时 PQ 左侧导轨的总电阻为 $R$，右侧导轨单位长度的电阻为 $R_{0}$。以 ef 为界，其左侧匀强磁场方向竖直向上，右侧匀强磁场水平向左，磁感应强度大小均为 $B$。在 $t =$ 0 时，一水平向左的拉力 $F$ 垂直作用于导轨的 bc 边上，使导轨由静止开始做匀加速直线运动，加速度为$a$ 。
\begin{enumerate}
	\item
	求回路中感应电动势及感应电流随时间变化的表达式；
	\item
	经过多少时间拉力 $F$ 达到最大值，拉力 $F$ 的最大值为多少？
	\item
	某一过程中回路产生的焦耳热为 $Q$，导轨克服摩擦力做功为 $W$，求导轨动能的增加量。
\end{enumerate}
\onepicture[h=4cm, align=right]{\includesvg{figs/33}}
%% answer: 
\jdanswer{
\begin{enumerate}
	\item
	$I = \frac{BLat}{R + R_{0}at^{2}}$

	\item
	当 $t = \sqrt {\frac{R}{aR_{0}}} $ 时外力 $F$ 取最大值，$F_{max} = Ma + \mu mg + \frac{(1 + \mu )B^{2}L^{2}}{2}\sqrt {\frac{a}{RR_{0}}} $。

	\item
	$\Delta E_{k} = \frac{Ma}{\mu mg}$（$W - \mu Q$）

\end{enumerate}
}
%% memo:
\memoanswer{
（1）感应电动势为 $E = BLv$

导轨做初速为零的匀加速运动

$v = at$，$s = \frac{1}{2}at^{2}$

感应电动势 $E = BLat$

回路中感应电流随时间变化的表达式

$I = \frac{{BLv}}{{{R_\text{总}}}} = \frac{{BLat}}{{R + 2{R_0} \times \frac{{a{t^2}}}{2}}} = \frac{{BLat}}{{R + {R_0}a{t^2}}}$

（2）导轨受外力 $F$，安培力$F_{A}$，摩擦力 $F_{f}$。其中

$F_{A} = BIL = \frac{{{B^2}{L^2}at}}{{R + {R_0}a{t^2}}}$

$F_{f} = \mu F_{N} = \mu$（$mg + BIL$）$=\mu$（$mg + \frac{{{B^2}{L^2}at}}{{R + {R_0}a{t^2}}}$）

由牛顿定律

$F - F_{A} - F_{f} = Ma$

$F = Ma + F_{A} + F_{f} = Ma + \mu mg +$（1 $+ \mu$）$\frac{{{B^2}{L^2}at}}{{R + {R_0}a{t^2}}}$

上式中当 $\frac{R}{t} = R_{0}at$，即 $t = \sqrt {\frac{R}{{a{R_0}}}} $ 时外力$ F$ 取最大值。

$F_{max} = Ma + \mu mg + \frac{{(1 + \mu ){B^2}{L^2}}}{2}\sqrt {\frac{a}{{R{R_0}}}} $

（3）设此过程中导轨运动距离为 $s$，由动能定理

$W_{\text{合}} = \Delta E_{k}$，$W_{\text{合}} = Mas$

由于摩擦力为 $F_{f} = \mu$（$mg + F_{A}$），所以摩擦力做功为

$W = \mu mgs + \mu W_{A} = \mu mgs + \mu Q$

$s = \frac{{W - \mu Q}}{{\mu mg}}$

$\Delta E_{k} = Mas = \frac{{Ma}}{{\mu mg}}$（$W - \mu Q$）
}

\end{enumerate}

\end{document}
